% GENERATED FILE. Edit canonical lessons, then run npm run generate:books.
% canonical-source-hash: fnv1a64:67928bfcb1d466bd
% canonical-lessons: MR-C401-indhan, MR-C401-petrol, MR-C401-tayar, MR-C401-pankcar, MR-C401-sampne

\chapter{Fuel and Tyres}
\label{ch:mr-fuel-and-tyres}
\hlchaptermodality{401}

\begin{chapteropening}
\textbf{By the end of this chapter:} \emph{I can talk about fuel, petrol, tyres and a puncture.}

Complete the last lesson of chapter 401: I can talk about fuel, petrol, tyres and a puncture.
\end{chapteropening}

\section[\texorpdfstring{\mr{इंधन} (indhan)}{इंधन (indhan)}]{\mr{इंधन} (indhan) — fuel}

\label{lesson:MR-C401-indhan}



\begin{quote}
Before the new one: say the Marathi for a truck, then the Marathi for to stuff in.
\end{quote}

\subsection*{You'll want to know: \mr{इंधन}}
\textbf{\mr{इंधन}} — \emph{indhan} — \textquotedblleft{}fuel\textquotedblright{}.

\begin{quote}
\mr{गाडीतलं} \mr{इंधन} \mr{कमी} \mr{झालं} \mr{आहे}.
\end{quote}

\begin{quote}
\emph{gāḍītlaṁ indhan kamī jhālaṁ āhe.}
\end{quote}

\begin{quote}
The fuel in the car has run low.
\end{quote}

Neuter.

Any fuel: wood, gas or petrol; \mr{इंधन} \mr{भरणे}, to fill up.

\begin{cousinweb}
\textbf{\mr{इंधन}} (\emph{indhan}) is a tatsama, taken from Sanskrit \emph{indhana}, \textquotedblleft{}kindling; fuel\textquotedblright{} (Monier-Williams).
\end{cousinweb}

\begin{grammarlens}[title={What you've built}]
One more word for getting around.
\end{grammarlens}

\subsection*{Your turn}
\begin{itemize}
\raggedright
  \item \emph{Say it:} \emph{indhan}
  \item \emph{Say it:} \emph{indhan}, once more
  \item \emph{Say it:} \emph{gāḍītlaṁ indhan kamī jhālaṁ āhe.} — The fuel in the car has run low.
  \item \emph{Recall:} say the Marathi for a truck, then the Marathi for to stuff in, then say \emph{indhan} again
\end{itemize}

\begin{tcolorbox}[breakable,colback=teal!4,colframe=teal!35!black,title={Before you move on}]
What does \mr{इंधन} mean? (\textquotedblleft{}Fuel\textquotedblright{}.) Say it once more.
\end{tcolorbox}
\section[\texorpdfstring{\mr{पेट्रोल} (peṭrol)}{पेट्रोल (peṭrol)}]{\mr{पेट्रोल} (peṭrol) — petrol}

\label{lesson:MR-C401-petrol}



\begin{quote}
Before the new one: say the Marathi for to stuff in, then the Marathi for fuel.
\end{quote}

\subsection*{You'll want to know: \mr{पेट्रोल}}
\textbf{\mr{पेट्रोल}} — \emph{peṭrol} — \textquotedblleft{}petrol\textquotedblright{}.

\begin{quote}
\mr{पुढच्या} \mr{गावात} \mr{पेट्रोल} \mr{मिळेल}.
\end{quote}

\begin{quote}
\emph{puḍhacyā gāvāt peṭrol miḷel.}
\end{quote}

\begin{quote}
We will get petrol in the next village.
\end{quote}

Neuter.

\mr{पेट्रोल} \mr{भरणे}, to fill up with petrol; \mr{मिळणे} takes the dative: \mr{मला} \mr{पेट्रोल} \mr{मिळालं}, I got petrol.

\begin{cousinweb}
\textbf{\mr{पेट्रोल}} (\emph{peṭrol}) is a loan from English \emph{petrol}.
\end{cousinweb}

\begin{grammarlens}[title={What you've built}]
One more word for getting around.
\end{grammarlens}

\subsection*{Your turn}
\begin{itemize}
\raggedright
  \item \emph{Say it:} \emph{peṭrol}
  \item \emph{Say it:} \emph{peṭrol}, once more
  \item \emph{Say it:} \emph{puḍhacyā gāvāt peṭrol miḷel.} — We will get petrol in the next village.
  \item \emph{Recall:} say the Marathi for to stuff in, then the Marathi for fuel, then say \emph{peṭrol} again
\end{itemize}

\begin{tcolorbox}[breakable,colback=teal!4,colframe=teal!35!black,title={Before you move on}]
What does \mr{पेट्रोल} mean? (\textquotedblleft{}Petrol\textquotedblright{}.) Say it once more.
\end{tcolorbox}
\section[\texorpdfstring{\mr{टायर} (ṭāyar)}{टायर (ṭāyar)}]{\mr{टायर} (ṭāyar) — a tyre}

\label{lesson:MR-C401-tayar}



\begin{quote}
Before the new one: say the Marathi for fuel, then the Marathi for petrol.
\end{quote}

\subsection*{You'll want to know: \mr{टायर}}
\textbf{\mr{टायर}} — \emph{ṭāyar} — \textquotedblleft{}a tyre\textquotedblright{}.

\begin{quote}
\mr{दुचाकीचा} \mr{मागचा} \mr{टायर} \mr{जुना} \mr{झाला} \mr{आहे}.
\end{quote}

\begin{quote}
\emph{ducākīcā māgcā ṭāyar junā jhālā āhe.}
\end{quote}

\begin{quote}
The back tyre of the two-wheeler has worn out.
\end{quote}

Masculine.

The plural is the same word: \mr{चार} \mr{टायर}.

\begin{cousinweb}
\textbf{\mr{टायर}} (\emph{ṭāyar}) is a loan from English \emph{tyre}.
\end{cousinweb}

\begin{grammarlens}[title={What you've built}]
One more word for getting around.
\end{grammarlens}

\subsection*{Your turn}
\begin{itemize}
\raggedright
  \item \emph{Say it:} \emph{ṭāyar}
  \item \emph{Say it:} \emph{ṭāyar}, once more
  \item \emph{Say it:} \emph{ducākīcā māgcā ṭāyar junā jhālā āhe.} — The back tyre of the two-wheeler has worn out.
  \item \emph{Recall:} say the Marathi for fuel, then the Marathi for petrol, then say \emph{ṭāyar} again
\end{itemize}

\begin{tcolorbox}[breakable,colback=teal!4,colframe=teal!35!black,title={Before you move on}]
What does \mr{टायर} mean? (\textquotedblleft{}A tyre\textquotedblright{}.) Say it once more.
\end{tcolorbox}
\section[\texorpdfstring{\mr{पंक्चर} (paṅkcar)}{पंक्चर (paṅkcar)}]{\mr{पंक्चर} (paṅkcar) — a puncture; a flat tyre}

\label{lesson:MR-C401-pankcar}



\begin{quote}
Before the new one: say the Marathi for petrol, then the Marathi for a tyre.
\end{quote}

\subsection*{You'll want to know: \mr{पंक्चर}}
\textbf{\mr{पंक्चर}} — \emph{paṅkcar} — \textquotedblleft{}a puncture; a flat tyre\textquotedblright{}.

\begin{quote}
\mr{वाटेत} \mr{सायकलचं} \mr{चाक} \mr{पंक्चर} \mr{झालं}.
\end{quote}

\begin{quote}
\emph{vāṭet sāykalcaṁ cāk paṅkcar jhālaṁ.}
\end{quote}

\begin{quote}
On the way the bicycle wheel got a puncture.
\end{quote}

Neuter.

\mr{पंक्चर} \mr{होणे}, to get a puncture; \mr{पंक्चर} \mr{काढणे}, to mend one.

\begin{cousinweb}
\textbf{\mr{पंक्चर}} (\emph{paṅkcar}) is a loan from English \emph{puncture}.
\end{cousinweb}

\begin{grammarlens}[title={What you've built}]
One more word for getting around.
\end{grammarlens}

\subsection*{Your turn}
\begin{itemize}
\raggedright
  \item \emph{Say it:} \emph{paṅkcar}
  \item \emph{Say it:} \emph{paṅkcar}, once more
  \item \emph{Say it:} \emph{vāṭet sāykalcaṁ cāk paṅkcar jhālaṁ.} — On the way the bicycle wheel got a puncture.
  \item \emph{Recall:} say the Marathi for petrol, then the Marathi for a tyre, then say \emph{paṅkcar} again
\end{itemize}

\begin{tcolorbox}[breakable,colback=teal!4,colframe=teal!35!black,title={Before you move on}]
What does \mr{पंक्चर} mean? (\textquotedblleft{}A puncture; a flat tyre\textquotedblright{}.) Say it once more.
\end{tcolorbox}
\section[\texorpdfstring{\mr{संपणे} (sampṇe)}{संपणे (sampṇe)}]{\mr{संपणे} (sampṇe) — to run out; to come to an end}

\label{lesson:MR-C401-sampne}



\begin{quote}
Before the new one: say the Marathi for a tyre, then the Marathi for a puncture.
\end{quote}

\subsection*{You'll want to know: \mr{संपणे}}
\textbf{\mr{संपणे}} — \emph{sampṇe} — \textquotedblleft{}to run out; to come to an end\textquotedblright{}.

\begin{quote}
\mr{रस्त्यात} \mr{पेट्रोल} \mr{संपलं}.
\end{quote}

\begin{quote}
\emph{rastyāt peṭrol samplaṁ.}
\end{quote}

\begin{quote}
The petrol ran out on the way.
\end{quote}

Past (he …): \mr{तो} \mr{संपला}, \emph{to samplā}.

What runs out is the subject: \mr{पैसे} \mr{संपले}, the money ran out; \mr{सुट्टी} \mr{संपली}, the holiday ended.

\begin{cousinweb}
\textbf{\mr{संपणे}} (\emph{sampṇe}) is a common Marathi verb, \textquotedblleft{}to be finished\textquotedblright{}; its origin is not settled.
\end{cousinweb}

\begin{grammarlens}[title={What you've built}]
One more word for getting around.
\end{grammarlens}

\subsection*{Your turn}
\begin{itemize}
\raggedright
  \item \emph{Say it:} \emph{sampṇe}
  \item \emph{Say it:} \emph{sampṇe}, once more
  \item \emph{Say it:} \emph{rastyāt peṭrol samplaṁ.} — The petrol ran out on the way.
  \item \emph{Recall:} say the Marathi for a tyre, then the Marathi for a puncture, then say \emph{sampṇe} again
\end{itemize}

\begin{tcolorbox}[breakable,colback=teal!4,colframe=teal!35!black,title={Before you move on}]
What does \mr{संपणे} mean? (\textquotedblleft{}To run out; to come to an end\textquotedblright{}.) Say it once more.
\end{tcolorbox}
